\documentclass[aspectratio=169,11pt]{beamer}
\usetheme{Madrid}
\usecolortheme{dolphin}
\setbeamertemplate{navigation symbols}{}
\setbeamertemplate{caption}[numbered]

\usepackage[utf8]{inputenc}
\usepackage[T1]{fontenc}
\usepackage[spanish,es-noshorthands]{babel}
\usepackage{amsmath,amssymb,amsthm,mathtools}
\usepackage{tikz}
\usepackage{pgfplots}
\pgfplotsset{compat=1.18}
\usepackage{booktabs}
\usepackage{array}

\title[Prob. y Estadística]{Clase 25: Conceptos de prueba de hipótesis}
\subtitle{Unidad 7: Prueba de hipótesis}
\institute[UNS]{Universidad Nacional del Sur \\ Departamento de Matemática}
\author{Probabilidad y Estadística (5765)}
\date{}

\begin{document}

\begin{frame}
\titlepage
\end{frame}

\begin{frame}{Contenidos}
\tableofcontents
\end{frame}

\section{Introducción}

\begin{frame}{De la estimación a la prueba de hipótesis}
En la Unidad 6 estudiamos cómo \textbf{estimar} un parámetro poblacional desconocido $\theta$ (puntual o por intervalo de confianza) a partir de una muestra aleatoria.

\begin{block}{Un problema distinto}
Muchas veces el objetivo no es estimar $\theta$, sino \textbf{decidir}, con la evidencia muestral disponible, si una afirmación (conjetura) sobre $\theta$ es sostenible o debe ser rechazada.
\end{block}

\begin{exampleblock}{Ejemplos de preguntas}
\begin{itemize}
\item ¿El contenido medio de las bolsas de un producto es realmente $500$ g, como indica el rótulo?
\item ¿Un nuevo proceso de fabricación reduce la proporción de piezas defectuosas?
\item ¿La variabilidad de un instrumento de medición supera la tolerancia admitida?
\item ¿El rendimiento académico es independiente del turno de cursada?
\end{itemize}
\end{exampleblock}

Estas preguntas se responden mediante una \textbf{prueba (o contraste) de hipótesis}, un procedimiento formal de decisión estadística.
\end{frame}

\section{Hipótesis estadísticas}

\begin{frame}{Hipótesis nula y alternativa}
\begin{definition}
Una \textbf{hipótesis estadística} es una afirmación (conjetura) acerca del valor de uno o más parámetros de una o varias poblaciones.
\end{definition}

\begin{block}{Hipótesis nula $H_0$}
Es la afirmación que se somete a prueba. Representa la situación de referencia, el status quo, o la ausencia de efecto/diferencia. Se formula siempre con una igualdad:
\[
H_0: \theta = \theta_0.
\]
\end{block}

\begin{block}{Hipótesis alternativa $H_1$ (o $H_a$)}
Es la afirmación contraria a $H_0$, la que se acepta como razonable si la evidencia muestral contradice a $H_0$. Puede ser:
\[
H_1: \theta \neq \theta_0 \quad \text{(bilateral)}, \qquad H_1:\theta>\theta_0 \ \text{ o } \ H_1:\theta<\theta_0 \quad \text{(unilateral)}.
\]
\end{block}
\end{frame}

\begin{frame}{Cómo se formulan $H_0$ y $H_1$}
\begin{alertblock}{Principio general}
$H_0$ y $H_1$ deben ser mutuamente excluyentes y, en la práctica usual de esta materia, exhaustivas respecto del parámetro. La afirmación que se quiere \emph{probar con evidencia} (aquello que el investigador sospecha o desea demostrar) suele plantearse como $H_1$, mientras que $H_0$ representa lo contrario o el criterio conservador.
\end{alertblock}

\begin{example}
Una fábrica afirma que el diámetro medio de sus rodamientos es $\mu_0 = 10$ mm. Un control de calidad sospecha que el proceso se desvió. Se plantea:
\[
H_0: \mu = 10 \quad \text{vs.} \quad H_1: \mu \neq 10.
\]
\end{example}

\begin{example}
Un nuevo fertilizante se prueba para ver si \textbf{aumenta} el rendimiento medio $\mu$ respecto del valor histórico $\mu_0=40$ qq/ha:
\[
H_0: \mu \le 40 \quad \text{vs.} \quad H_1: \mu > 40.
\]
En la práctica, para el cálculo se trabaja con $H_0:\mu=40$, el punto frontera más desfavorable.
\end{example}
\end{frame}

\section{Tipos de error}

\begin{frame}{La decisión y sus errores posibles}
Una prueba de hipótesis conduce a una de dos decisiones: \textbf{rechazar} $H_0$ o \textbf{no rechazar} $H_0$. Como la decisión se basa en una muestra, siempre existe la posibilidad de equivocarse.

\begin{table}[]
\centering
\begin{tabular}{@{}lcc@{}}
\toprule
 & $H_0$ verdadera & $H_0$ falsa \\
\midrule
\textbf{No rechazar $H_0$} & Decisión correcta ($1-\alpha$) & Error tipo II ($\beta$) \\
\textbf{Rechazar $H_0$} & Error tipo I ($\alpha$) & Decisión correcta ($1-\beta$) \\
\bottomrule
\end{tabular}
\end{table}

\begin{block}{Error de tipo I}
Rechazar $H_0$ cuando en realidad es verdadera. Su probabilidad se llama \textbf{nivel de significación}:
\[
\alpha = P(\text{Rechazar } H_0 \mid H_0 \text{ verdadera}).
\]
\end{block}

\begin{block}{Error de tipo II}
No rechazar $H_0$ cuando en realidad es falsa. Su probabilidad se denota:
\[
\beta = P(\text{No rechazar } H_0 \mid H_0 \text{ falsa}).
\]
\end{block}
\end{frame}

\begin{frame}{Relación entre $\alpha$ y $\beta$}
\begin{itemize}
\item $\alpha$ se fija de antemano por el investigador (valores usuales: $0.10$, $0.05$, $0.01$). Es la probabilidad máxima que se está dispuesto a tolerar de rechazar erróneamente una $H_0$ verdadera.
\item Para un tamaño de muestra $n$ fijo, $\alpha$ y $\beta$ están inversamente relacionados: si se achica la región de rechazo para disminuir $\alpha$, aumenta $\beta$, y viceversa.
\item La \textbf{única} forma de disminuir simultáneamente $\alpha$ y $\beta$ es aumentar el tamaño muestral $n$.
\end{itemize}

\begin{alertblock}{¿Por qué se prioriza $\alpha$?}
$H_0$ suele representar la afirmación "conservadora" o la que produce consecuencias más costosas si se rechaza equivocadamente (por ejemplo, retirar un lote en buen estado, acusar de fraude a un proceso correcto). Por eso se controla directamente $\alpha$, y $\beta$ se analiza mediante la potencia.
\end{alertblock}
\end{frame}

\section{Potencia de la prueba}

\begin{frame}{Potencia de una prueba}
\begin{definition}
La \textbf{potencia} de una prueba de hipótesis es la probabilidad de rechazar $H_0$ cuando efectivamente es falsa:
\[
\text{Potencia} = 1-\beta = P(\text{Rechazar } H_0 \mid H_0 \text{ falsa}).
\]
\end{definition}

\begin{block}{Interpretación}
La potencia mide la capacidad de la prueba para \textbf{detectar} que $H_0$ es falsa cuando el verdadero parámetro $\theta$ se aparta de $\theta_0$ en una magnitud dada. Es función del valor real del parámetro, del tamaño muestral $n$ y del nivel $\alpha$ elegido.
\end{block}

\begin{itemize}
\item Cuanto más se aleja el valor real de $\theta_0$, mayor es la potencia (es "más fácil" detectar diferencias grandes).
\item A mayor $n$, mayor potencia para una misma diferencia a detectar.
\item A mayor $\alpha$ (región de rechazo más amplia), mayor potencia, pero también mayor riesgo de error tipo I.
\end{itemize}

\begin{alertblock}{Una prueba ``buena''}
Entre pruebas con el mismo $\alpha$, se prefiere la de mayor potencia (menor $\beta$) para detectar el efecto de interés.
\end{alertblock}
\end{frame}

\section{Estadístico de prueba y región crítica}

\begin{frame}{Estadístico de prueba}
\begin{definition}
El \textbf{estadístico de prueba} (o de contraste) es una función de la muestra aleatoria, $T=T(X_1,\dots,X_n)$, cuya distribución de probabilidad \textbf{bajo el supuesto de que $H_0$ es verdadera} es conocida.
\end{definition}

\begin{block}{Rol del estadístico}
Resume la información muestral relevante para el parámetro en cuestión y permite comparar lo observado con lo que se esperaría si $H_0$ fuera cierta. Valores del estadístico ``poco compatibles'' con $H_0$ llevan a rechazarla.
\end{block}

\begin{example}
Para probar $H_0:\mu=\mu_0$ con varianza poblacional conocida $\sigma^2$, el estadístico de prueba es
\[
Z = \frac{\overline{X}-\mu_0}{\sigma/\sqrt{n}} \ \sim\ N(0,1) \quad \text{bajo } H_0.
\]
\end{example}
\end{frame}

\begin{frame}{Región crítica y región de aceptación}
\begin{definition}
El espacio de todos los valores posibles del estadístico de prueba se divide en dos subconjuntos:
\begin{itemize}
\item \textbf{Región crítica o de rechazo} $RR$: valores del estadístico que llevan a rechazar $H_0$.
\item \textbf{Región de aceptación (no rechazo)} $RA$: el complemento, valores compatibles con $H_0$.
\end{itemize}
\end{definition}

\begin{block}{Construcción de la región crítica}
Se determina de modo que, bajo $H_0$, la probabilidad de que el estadístico caiga en $RR$ sea exactamente $\alpha$ (el nivel de significación fijado). La forma de $RR$ (una o dos colas) depende de $H_1$.
\end{block}

\begin{itemize}
\item $H_1:\theta\neq\theta_0$ (bilateral) $\Rightarrow$ $RR$ en \textbf{ambas colas} de la distribución.
\item $H_1:\theta>\theta_0$ $\Rightarrow$ $RR$ en la \textbf{cola derecha}.
\item $H_1:\theta<\theta_0$ $\Rightarrow$ $RR$ en la \textbf{cola izquierda}.
\end{itemize}
\end{frame}

\begin{frame}{Ilustración: prueba bilateral vs. unilateral}
\centering
\begin{tikzpicture}[scale=0.85]
\begin{axis}[
  width=6.2cm, height=3.6cm,
  axis lines=left, ytick=\empty,
  xtick={-1.96,0,1.96}, xticklabels={$-z_{\alpha/2}$,$0$,$z_{\alpha/2}$},
  title={\footnotesize Bilateral: $H_1:\mu\neq\mu_0$},
  ymin=0, ymax=0.45,
]
\addplot[domain=-3.5:3.5, samples=100, thick, blue] {exp(-x^2/2)/sqrt(2*pi)};
\addplot[domain=-3.5:-1.96, samples=40, fill=red!40, draw=none] {exp(-x^2/2)/sqrt(2*pi)} \closedcycle;
\addplot[domain=1.96:3.5, samples=40, fill=red!40, draw=none] {exp(-x^2/2)/sqrt(2*pi)} \closedcycle;
\end{axis}
\end{tikzpicture}
\hspace{0.3cm}
\begin{tikzpicture}[scale=0.85]
\begin{axis}[
  width=6.2cm, height=3.6cm,
  axis lines=left, ytick=\empty,
  xtick={0,1.645}, xticklabels={$0$,$z_{\alpha}$},
  title={\footnotesize Unilateral: $H_1:\mu>\mu_0$},
  ymin=0, ymax=0.45,
]
\addplot[domain=-3.5:3.5, samples=100, thick, blue] {exp(-x^2/2)/sqrt(2*pi)};
\addplot[domain=1.645:3.5, samples=40, fill=red!40, draw=none] {exp(-x^2/2)/sqrt(2*pi)} \closedcycle;
\end{axis}
\end{tikzpicture}

\vspace{0.3cm}
Las regiones sombreadas (rojas) son las regiones críticas, de área total $\alpha$.
\end{frame}

\section{Valor p}

\begin{frame}{El valor p (p-value)}
\begin{definition}
El \textbf{valor p} es la probabilidad, calculada bajo el supuesto de que $H_0$ es verdadera, de obtener un valor del estadístico de prueba tan extremo o más extremo (en el sentido indicado por $H_1$) que el efectivamente observado en la muestra.
\end{definition}

\begin{block}{Regla de decisión con el valor p}
\[
\text{Si } p\text{-valor} \le \alpha \ \Rightarrow \ \text{se rechaza } H_0.
\]
\[
\text{Si } p\text{-valor} > \alpha \ \Rightarrow \ \text{no se rechaza } H_0.
\]
\end{block}

\begin{alertblock}{Interpretación}
El valor p \textbf{no} es la probabilidad de que $H_0$ sea verdadera. Es una medida de cuán compatible es la evidencia muestral con $H_0$: cuanto menor el valor p, mayor la evidencia \emph{en contra} de $H_0$.
\end{alertblock}

Es equivalente decidir comparando el estadístico con el valor crítico o comparando el valor p con $\alpha$: ambos métodos siempre coinciden.
\end{frame}

\section{Pasos para realizar una prueba}

\begin{frame}{Procedimiento general de una prueba de hipótesis}
\begin{enumerate}
\item Formular $H_0$ y $H_1$ en términos del parámetro poblacional de interés.
\item Fijar el nivel de significación $\alpha$.
\item Elegir el estadístico de prueba adecuado y establecer su distribución bajo $H_0$.
\item Determinar la región crítica (valor(es) crítico(s)) de acuerdo con $\alpha$ y con la forma de $H_1$ (bilateral o unilateral), \emph{o} calcular el valor p.
\item Calcular el valor del estadístico de prueba a partir de los datos muestrales.
\item \textbf{Decidir}: rechazar $H_0$ si el estadístico cae en la región crítica (o si el valor p $\le \alpha$); no rechazar $H_0$ en caso contrario.
\item \textbf{Concluir} en el contexto del problema, con lenguaje no técnico.
\end{enumerate}
\end{frame}

\begin{frame}{Observaciones importantes}
\begin{itemize}
\item ``No rechazar $H_0$'' \textbf{no equivale} a ``aceptar $H_0$ como verdadera''. Significa que la evidencia muestral no es suficiente para rechazarla al nivel $\alpha$ elegido.
\item El nivel $\alpha$ se elige \textbf{antes} de observar los datos, nunca después.
\item Cambiar $\alpha$, $n$ o la forma de $H_1$ \emph{después} de ver los resultados invalida el procedimiento.
\item Estos mismos principios (estadístico, distribución bajo $H_0$, región crítica, decisión) se aplican a todas las pruebas específicas que veremos en las próximas clases: para medias, varianzas, proporciones, bondad de ajuste, independencia y homogeneidad.
\end{itemize}
\end{frame}

\section{Relación con los intervalos de confianza}

\begin{frame}{Prueba de hipótesis vs. intervalo de confianza}
\begin{block}{Dualidad fundamental}
Para una prueba bilateral $H_0:\theta=\theta_0$ vs. $H_1:\theta\neq\theta_0$ con nivel $\alpha$, y un intervalo de confianza para $\theta$ con nivel de confianza $1-\alpha$ construido con el mismo estadístico:
\[
\text{Se rechaza } H_0 \text{ al nivel } \alpha \quad \Longleftrightarrow \quad \theta_0 \notin IC_{1-\alpha}(\theta).
\]
\end{block}

\begin{example}
Si un intervalo de confianza del $95\%$ para $\mu$ resulta $IC = (48.3;\,51.7)$ y se desea probar $H_0:\mu=50$ vs. $H_1:\mu\neq50$ con $\alpha=0.05$: como $50 \in (48.3;51.7)$, \textbf{no se rechaza} $H_0$.
\end{example}

\begin{alertblock}{Utilidad práctica}
Un único intervalo de confianza permite decidir simultáneamente sobre \emph{cualquier} valor hipotético $\theta_0$ bilateral, mientras que cada prueba de hipótesis evalúa un único valor $\theta_0$ a la vez. Esta equivalencia es válida para pruebas bilaterales; para unilaterales se usa la correspondencia con intervalos de confianza de una cola.
\end{alertblock}
\end{frame}

\section{Resumen}

\begin{frame}{Resumen de la clase}
\begin{itemize}
\item Una \textbf{prueba de hipótesis} decide, con evidencia muestral, entre $H_0$ (hipótesis nula) y $H_1$ (alternativa).
\item \textbf{Error tipo I} ($\alpha$): rechazar $H_0$ siendo verdadera. \textbf{Error tipo II} ($\beta$): no rechazar $H_0$ siendo falsa.
\item La \textbf{potencia} $1-\beta$ mide la capacidad de detectar que $H_0$ es falsa.
\item El \textbf{estadístico de prueba} tiene distribución conocida bajo $H_0$; su valor se compara con la \textbf{región crítica}, determinada por $\alpha$ y la forma de $H_1$.
\item El \textbf{valor p} resume la evidencia contra $H_0$: se rechaza si $p\text{-valor}\le\alpha$.
\item Procedimiento: formular hipótesis $\to$ fijar $\alpha$ $\to$ estadístico y su distribución $\to$ región crítica $\to$ calcular $\to$ decidir $\to$ concluir.
\item Existe una \textbf{dualidad} entre pruebas bilaterales e intervalos de confianza.
\item Próximas clases: pruebas específicas para medias y varianzas de una y dos poblaciones, y pruebas basadas en $\chi^2$.
\end{itemize}
\end{frame}

\end{document}
